\section{Related Work}\label{sec:related}

\para{Stochastic neurons and stochastic activations.} The most direct precedent for our question replaces a deterministic hidden activation with a sample. \citet{bengio2013stochastic} draw a binary value $h_i \sim \mathrm{Bernoulli}(\sigma(a_i))$ at each hidden unit and introduce the straight-through estimator to pass gradients through the non-differentiable sample; this is exactly our \bernoullist{} mode. \citet{lee2019probact} push the idea to every layer with ProbAct, a stochastic activation function whose output is sampled from a per-element (mean, variance) distribution with a learnable variance, reporting generalization gains. These methods establish that hidden-layer sampling is trainable and can help, but they fix the operator (Bernoulli, or a single Gaussian-like form) rather than comparing distribution types. We place the Bernoulli stochastic neuron on the same axis as continuous and categorical sampling and measure where each one falls.

\para{Reparameterization and differentiable sampling.} Sampling at a hidden layer blocks gradients unless the sample is made differentiable. For continuous activations, the reparameterization trick \citep{kingma2014vae,rezende2014stochastic} writes $z = \mu + \sigma \odot \epsilon$ with parameter-free noise $\epsilon$, giving a low-variance pathwise gradient; this is the machinery behind our \gaussreparam{} mode. For categorical distributions, Gumbel-Softmax and the equivalent Concrete distribution \citep{jang2017gumbel,maddison2017concrete} relax a one-hot sample of $\softmax(x)$ through a temperature-$\tau$ continuous approximation, recovering a hard sample as $\tau \to 0$. This is the most direct match to our motivating question: sampling a hidden layer the way we already sample the output softmax. We use it for \gumbelsoft{} (soft, $\tau>0$) and \gumbelst{} (hard one-hot via straight-through). Our contribution is to put the continuous and categorical reparameterizations on the same axis and measure the gap: the Gaussian path costs essentially nothing, whereas the categorical path carries the highest gradient variance and the largest accuracy penalty.

\para{Discrete representation learning.} VQ-VAE \citep{oord2017vqvae} places a categorical codebook at an intermediate layer and trains it with a straight-through gradient, demonstrating that discrete hidden representations are learnable. Crucially, it warns that a soft relaxation of the codebook can be inverted or ignored by a sufficiently strong downstream stack, and that the discrete hidden code can collapse to a few entries. Our \cifar{} result instantiates exactly this caution: when the categorical sample collapses a hidden layer to roughly one active unit, the harder task cannot recover and \gumbelst{} falls to chance. Where VQ-VAE engineers around collapse to make discrete codes useful, we hold the architecture fixed and report collapse as a measured consequence of categorical hidden sampling.

\para{Stochasticity as regularization, Bayesian inference, and robustness.} A large body of work injects noise into a network for reasons other than studying activation sampling per se. Dropout and MC-dropout \citep{srivastava2014dropout,gal2016dropout} multiply activations by a Bernoulli mask, the latter reinterpreting test-time sampling as approximate Bayesian inference; variational dropout \citep{kingma2015variational} learns per-unit noise via the local reparameterization trick. Bayes-by-backprop \citep{blundell2015weight} instead samples the \emph{weights} to capture parameter uncertainty --- a deliberate contrast, since we sample activations and leave weights deterministic. Noisy activation functions \citep{gulcehre2016noisy} add noise to ease optimization of saturating nonlinearities, and the deep variational information bottleneck \citep{alemi2017deepvib,tishby2015deep} samples a stochastic latent to compress nuisance information. Stochastic activation pruning \citep{dhillon2018sap} samples a sparse activation mask for adversarial robustness, though those robustness claims were later weakened under stronger adversarial evaluation. We adopt \dropout{} and \gaussnoise{} as baselines from this lineage and quantify the regularization effect each operator produces, measuring calibration with the expected calibration error (\ece{}) of \citet{guo2017calibration} alongside accuracy and the generalization gap.

\para{Our position.} Across all of this work, stochasticity is introduced either for a \emph{purpose} --- generation, compression, uncertainty estimation, or robustness --- or at a \emph{fixed place} --- a latent bottleneck, a codebook, or every layer simultaneously. None isolates intermediate-layer activation sampling as the single variable of study under one protocol that sweeps distribution type, sampling amount, layer position, and train-versus-test usage while holding everything else constant. That is exactly what we do. By placing one \samplinglayer{} after a hidden ReLU and swapping only its sampling distribution, we make the analogy to output-softmax sampling explicit and ask what it costs to apply the same operation one layer earlier.
